
\section{The \texorpdfstring{$\e$}{e} channel: Abel normalization and the pole barrier}
\label{app:abel_pole_barrier}

This appendix records the minimal analytic stability viewpoint used in Section~\ref{subsubsec:e_channel}.
The golden-mean grammar defines a shift of finite type with adjacency matrix
\[
A=\begin{pmatrix}
1 & 1\\
1 & 0
\end{pmatrix},
\]
whose Artin--Mazur zeta function is
\[
\zeta(z)=\frac{1}{\det(I-zA)}=\frac{1}{1-z-z^2}
\]
\cite{ArtinMazur1965,ParryPollicott1990}.

\subsection{Spectral normalization}
The spectral radius is $\rho(A)=\varphi$, the golden ratio.
Introduce the normalized variable $z=r/\rho(A)=r/\varphi$.
Then the normalized zeta is
\[
\tilde{\zeta}(r):=\zeta(r/\varphi)=\frac{1}{(1-r)\left(1+\varphi^{-2}r\right)}.
\]
Therefore $\tilde{\zeta}$ is holomorphic on $|r|<1$, and its dominant pole lies on the unit-circle boundary at $r=1$.
This motivates a pole-barrier criterion for stability: admissible dynamics exhibit a barrier that prevents the leading singularity from entering the open unit disc under normalized scaling.

\subsection{Weighted extensions and pressure}
A weighted extension modifies the transfer operator without changing the forbidden grammar, yielding a family of matrices (or operators) $A_\beta$ with spectral radius $\rho(A_\beta)$.
In thermodynamic formalism, the pressure is $\log\rho(A_\beta)$ \cite{Ruelle1978}.
This provides a rigorous interface between analytic stability and a pressure-like quantity that can be interpreted as an information potential at higher resolution.



